\chapter{仮想環境とは何か}

\section{なぜ仮想環境が必要か}

Python でプログラムを書くとき、\cmd{numpy} や \cmd{matplotlib} などの
外部ライブラリ（パッケージ）を使います。
これらには\textbf{バージョン}があり、プロジェクトによって必要なバージョンが異なります。

\subsection{バージョン衝突の問題}

次のようなシナリオを考えてみましょう。

\begin{center}
\begin{tabularx}{\linewidth}{lXX}
  \toprule
  & B4課題（NMR ラボ） & 別の研究プログラム \\
  \midrule
  numpy & 1.26 が必要 & 2.0 が必要 \\
  \bottomrule
\end{tabularx}
\end{center}

同じ Python 環境に両方をインストールしようとすると、
片方が上書きされてどちらかが動かなくなります。

\begin{center}
\begin{tabularx}{\linewidth}{lX}
  \toprule
  状況 & 結果 \\
  \midrule
  numpy 1.26 をインストール & 別の研究プログラムが動かない \\
  numpy 2.0 にアップグレード & B4課題が動かなくなる \\
  \bottomrule
\end{tabularx}
\end{center}

これが\textbf{依存関係の衝突}です。

\subsection{グローバル環境を汚さないために}

\cmd{pip install} を繰り返すと、システム全体の Python 環境に
パッケージがどんどん蓄積されます。

\begin{itemize}
  \item 使わなくなったパッケージが残り続ける
  \item バージョンアップで他のプログラムが突然壊れる
  \item 「先輩の PC では動いたのに自分の PC では動かない」が起きやすい
\end{itemize}

\section{仮想環境の仕組み}

\textbf{仮想環境（Virtual Environment）}はプロジェクトごとに
独立した Python の箱を作る仕組みです。

\begin{center}
\fbox{\parbox{0.85\textwidth}{
\textbf{仮想環境のイメージ}\\[0.5em]
システム全体\\
\quad└ Ubuntu のシステム Python（素の状態）\\
\quad\quad ├ B4課題用 仮想環境（\cmd{.venv}）\\
\quad\quad │ \quad numpy 1.26, bm3d 4.0, jupyterlab 4.0 ...\\
\quad\quad └ 別プロジェクト用 仮想環境（\cmd{.venv}）\\
\quad\quad \quad \quad numpy 2.0, tensorflow ...
}}
\end{center}

各仮想環境は\textbf{完全に独立}しており、一方を変更しても他方には影響しません。

\section{仮想環境の実体}

仮想環境はただのフォルダです。
プロジェクトフォルダの中に \cmd{.venv/} というフォルダが作られ、
その中にライブラリの実体が格納されます。

\begin{wslbox}[.venv の中身のイメージ]
\begin{lstlisting}[style=bash]
2026年度_B4課題引継ぎ/
  ├── pyproject.toml   # どのライブラリが必要か書いてある
  ├── uv.lock          # バージョンを固定するファイル
  ├── .venv/           # 仮想環境（uv sync で自動生成）
  │     ├── bin/
  │     └── lib/python3.11/site-packages/
  │           ├── numpy/
  │           ├── matplotlib/
  │           └── ...
  ├── No.1_FID/
  └── ...
\end{lstlisting}
\end{wslbox}

\begin{tipbox}
\cmd{.venv/} の中身は自動生成されるため、自分で触る必要はありません。
\cmd{uv sync} を実行するだけで作られます。
\end{tipbox}

\section{仮想環境を使うメリット}

\begin{description}
  \item[\textbf{依存関係の分離}]
    プロジェクトごとに必要なパッケージだけをインストールできる
  \item[\textbf{再現性}]
    全員が同じバージョンのライブラリを使えるので
    「自分の環境では動くのに\ldots\ldots」という問題が起きない
  \item[\textbf{削除が簡単}]
    不要になったら \cmd{.venv/} フォルダを削除するだけ
  \item[\textbf{システムを汚さない}]
    Ubuntu のシステム Python は変更しないので OS が安定する
\end{description}
